\chapter{API reference}\label{app:api}

This appendix lists every public class and function of \code{fuzzyclust}, sub-package by sub-package, with the first sentence of its docstring. The lists are generated from the \code{\_\_all\_\_} attribute of each module by \code{book/build\_api\_appendix.py}, so they are complete by construction; the full signatures, parameters, and attributes are in the online documentation and in the docstrings (\code{help(fuzzyclust.cluster.GK)}). The shared constructor parameters, fitted attributes, and methods of all clustering estimators are described once in \cref{sec:found:api}. The notation of the docstrings is that of \cref{tab:found:notation}.

{\footnotesize
% Generated by book/build_api_appendix.py -- do not edit by hand.

\section[\code{fuzzyclust.base}]{\code{fuzzyclust.base}: Base class}
\begin{longtable}{@{}>{\raggedright\arraybackslash}p{0.39\textwidth} l >{\raggedright\arraybackslash}p{0.43\textwidth}@{}}\toprule
Name & Kind & Description \\\midrule\endfirsthead
\toprule Name & Kind & Description \\\midrule\endhead
\bottomrule\endfoot
\code{BaseFuzzyClusterer} & class & Abstract base of the prototype-based clustering algorithms \\
\code{EPS} & constant & Small constant added to squared distances before negative powers ($10^{-10}$) \\
\end{longtable}

\section[\code{fuzzyclust.cluster}]{\code{fuzzyclust.cluster}: Partitioning algorithms}
\begin{longtable}{@{}>{\raggedright\arraybackslash}p{0.39\textwidth} l >{\raggedright\arraybackslash}p{0.43\textwidth}@{}}\toprule
Name & Kind & Description \\\midrule\endfirsthead
\toprule Name & Kind & Description \\\midrule\endhead
\bottomrule\endfoot
\code{KMeans} & class & Hard c-means (k-means) clustering with the fuzzyclust interface \\
\code{KMedoids} & class & Hard k-medoids clustering (prototypes are data points) \\
\code{FCM} & class & Fuzzy c-means clustering (Dunn, Bezdek) \\
\code{GK} & class & Gustafson--Kessel (GK) fuzzy clustering \\
\code{GG} & class & Gath--Geva (GG) fuzzy clustering \\
\code{RegularizedFCM} & class & Fuzzy c-means by entropy or quadratic regularization, with optional deterministic annealing \\
\code{PolynomialFuzzifierFCM} & class & Fuzzy c-means with the polynomial fuzzifier of Klawonn and H\"{o}ppner \\
\code{entropy\_memberships} & function & Memberships of the entropy-regularized c-means, $\mu_{ik}\propto\exp(-D^2_{ik}/\lambda)$ \\
\code{quadratic\_memberships} & function & Memberships of the quadratically regularized c-means (Miyamoto and Umayahara) \\
\code{polynomial\_memberships} & function & Memberships of fuzzy c-means with the polynomial fuzzifier (Klawonn and H\"{o}ppner) \\
\code{polynomial\_fuzzifier} & function & Polynomial fuzzifier $g(\mu)=\frac{1-\beta}{1+\beta}\mu^2+\frac{2\beta}{1+\beta}\mu$ of Klawonn and H\"{o}ppner \\
\code{SubtractiveClustering} & class & Subtractive clustering (Chiu, 1994): cluster centers from the potential of the data points \\
\code{MountainClustering} & class & Mountain method (Yager and Filev, 1994): cluster centers from a mountain function on a grid \\
\code{CompatibleClusterMerging} & class & Compatible cluster merging: find the number of clusters by merging from a large start \\
\code{cluster\_similarity} & function & Pairwise similarity of fuzzy clusters from their membership columns (fuzzy Jaccard index) \\
\code{geometric\_compatibility} & function & Compatibility of cluster pairs from centers and covariances (Krishnapuram and Freg) \\
\code{reference\_fuzziness} & function & Rescale an FCM-type partition obtained with exponent m to the one exponent m\_ref would give \\
\end{longtable}

\section[\code{fuzzyclust.validity}]{\code{fuzzyclust.validity}: Validity indices}
\begin{longtable}{@{}>{\raggedright\arraybackslash}p{0.39\textwidth} l >{\raggedright\arraybackslash}p{0.43\textwidth}@{}}\toprule
Name & Kind & Description \\\midrule\endfirsthead
\toprule Name & Kind & Description \\\midrule\endhead
\bottomrule\endfoot
\code{partition\_coefficient} & function & Partition coefficient (Bezdek) \\
\code{classification\_entropy} & function & Classification (partition) entropy (Bezdek) \\
\code{modified\_partition\_coefficient} & function & Modified partition coefficient (Dav\'{e}, 1996) \\
\code{partition\_index} & function & Partition index SC (Bensaid et al., 1996) \\
\code{separation\_index} & function & Separation index S: Xie--Beni in its original form, with squared memberships (Xie and Beni, 1991) \\
\code{xie\_beni} & function & Xie--Beni index (Xie and Beni, 1991) \\
\code{fukuyama\_sugeno} & function & Fukuyama--Sugeno index (1989) \\
\code{dunn\_index} & function & Dunn's index (1974) on the hardened partition \\
\code{alternative\_dunn\_index} & function & Alternative Dunn index (ADI) of the MATLAB toolbox \\
\code{silhouette} & function & Mean silhouette coefficient (Rousseeuw, 1987) of the hardened partition \\
\code{fuzzy\_hypervolume} & function & Fuzzy hypervolume $FHV = \sum_i \sqrt{\det\mathbf{F}_i}$ (Gath and Geva, 1989) \\
\code{partition\_density} & function & Partition density $PD = \sum_i S_i / FHV$ (Gath and Geva, 1989) \\
\code{average\_partition\_density} & function & Average partition density (Gath and Geva) \\
\code{pbmf\_index} & function & PBM-F index (Pakhira, Bandyopadhyay and Maulik, 2004) \\
\code{kwon\_index} & function & Kwon's index (1998) \\
\code{pcaes\_index} & function & Partition coefficient and exponential separation, PCAES (Wu and Yang, 2005) \\
\code{fuzzy\_rand\_index} & function & Fuzzy Rand index between two fuzzy partitions of the same samples \\
\code{fuzzy\_adjusted\_rand\_index} & function & Fuzzy adjusted Rand index (chance-corrected), from the fuzzy pair counts of fuzzy\_rand\_index \\
\code{fuzzy\_jaccard\_index} & function & Fuzzy Jaccard index $a/(a+b+c)$ from the fuzzy pair counts of fuzzy\_rand\_index \\
\code{stability\_selection} & function & Stability of the partition against resampling, for every number of clusters \\
\code{validity\_indices} & function & Evaluate several validity indices at once \\
\code{scan\_n\_clusters} & function & Fit the estimator for several numbers of clusters and tabulate validity indices \\
\code{suggest\_n\_clusters} & function & Optimal number of clusters suggested by every index of a scan table \\
\code{INDEX\_DIRECTION} & constant & For every validity index, whether its optimum is the maximum or the minimum \\
\end{longtable}

\section[\code{fuzzyclust.extensions}]{\code{fuzzyclust.extensions}: Robust and shape-oriented variants}
\begin{longtable}{@{}>{\raggedright\arraybackslash}p{0.39\textwidth} l >{\raggedright\arraybackslash}p{0.43\textwidth}@{}}\toprule
Name & Kind & Description \\\midrule\endfirsthead
\toprule Name & Kind & Description \\\midrule\endhead
\bottomrule\endfoot
\code{PCM} & class & Possibilistic c-means clustering (Krishnapuram and Keller) \\
\code{PFCM} & class & Possibilistic fuzzy c-means clustering (Pal, Pal, Keller and Bezdek) \\
\code{NoiseClustering} & class & Noise clustering: fuzzy c-means with a virtual noise cluster (Dav\'{e}) \\
\code{KernelFCM} & class & Kernel fuzzy c-means: FCM in the feature space of a Mercer kernel \\
\code{rbf\_kernel} & function & Gaussian (RBF) kernel matrix $K_{kl}=\exp(-\gamma\|\mathbf{x}_k-\mathbf{y}_l\|^2)$ \\
\code{KERNELS} & constant & Names of the kernels accepted by \code{KernelFCM} \\
\code{FuzzyCVarieties} & class & Fuzzy c-varieties (FCV) clustering: linear-variety prototypes (Bezdek et al.) \\
\code{FuzzyCShells} & class & Fuzzy c-(spherical) shells: circles / hyperspheres as prototypes \\
\code{VolumePrototypeFCM} & class & Fuzzy c-means with volume prototypes and optional adaptive cluster merging (Kaymak and Setnes, 2002) \\
\code{SCAD} & class & Simultaneous clustering and attribute discrimination (Frigui and Nasraoui, 2004) \\
\code{FeatureWeightedFCM} & class & Fuzzy c-means with global attribute weights (after Huang et al., 2005) \\
\code{plot\_shells} & function & Scatter the data and draw the fitted circles of a FuzzyCShells model \\
\code{plot\_varieties} & function & Scatter the data and draw the fitted lines of a FuzzyCVarieties model (n\_dims=1) \\
\code{plot\_typicality} & function & Scatter plot colored by the maximal typicality $\max_i t_{ik}$ of every sample \\
\code{make\_circles} & function & Intersecting circles (shell clusters) with optional arcs and outliers \\
\code{make\_close\_blobs} & function & Two Gaussian clusters of different size and density close to each other \\
\code{make\_crossing\_lines} & function & Line segments that cross at a common point \\
\code{make\_collinear\_segments} & function & Two collinear segments separated by a gap, optionally with a third, oblique segment \\
\end{longtable}

\section[\code{fuzzyclust.graph}]{\code{fuzzyclust.graph}: Graph-based clustering}
\begin{longtable}{@{}>{\raggedright\arraybackslash}p{0.39\textwidth} l >{\raggedright\arraybackslash}p{0.43\textwidth}@{}}\toprule
Name & Kind & Description \\\midrule\endfirsthead
\toprule Name & Kind & Description \\\midrule\endhead
\bottomrule\endfoot
\code{MSTClustering} & class & Hybrid clustering with a minimal spanning tree and Gath--Geva models \\
\code{GeodesicFCM} & class & Geodesic-distance-based fuzzy c-medoid clustering (Feil and Abonyi, Algorithm II) \\
\code{IsomapFCM} & class & Fuzzy c-means on the Isomap projection (Feil and Abonyi, Algorithm I) \\
\code{minimum\_spanning\_tree} & function & Euclidean minimal spanning tree of the samples \\
\code{fuzzy\_cluster\_similarity} & function & Pairwise similarity of fuzzy clusters seen as fuzzy sets over the samples \\
\code{neighbors\_graph} & function & Symmetric $k$-nearest-neighbor (or $\varepsilon$-radius) graph with Euclidean weights \\
\code{geodesic\_distances} & function & Graph (geodesic) distances between all pairs of samples, Isomap style \\
\code{isomap\_embedding} & function & Classical multidimensional scaling of a (geodesic) distance matrix (the Isomap projection) \\
\code{plot\_mst} & function & Draw the minimal spanning tree of a fitted MSTClustering \\
\code{plot\_neighbors\_graph} & function & Draw a neighborhood graph (sparse adjacency) over the data \\
\code{plot\_geodesic\_path} & function & Highlight the shortest graph path between two samples and the Euclidean chord \\
\code{plot\_similarity\_dendrogram} & function & Dendrogram of fuzzy clusters built from their similarity matrix \\
\code{make\_chained\_blobs} & function & Two Gaussian blobs connected by a sparse chain of points (the \emph{chaining effect}) \\
\code{make\_density\_blobs} & function & A diffuse cluster next to two tight ones (very different densities) \\
\code{make\_s\_curve} & function & Points on the three-dimensional S-curve, a two-dimensional manifold (wraps scikit-learn) \\
\end{longtable}

\section[\code{fuzzyclust.viz}]{\code{fuzzyclust.viz}: Visualization}
\begin{longtable}{@{}>{\raggedright\arraybackslash}p{0.39\textwidth} l >{\raggedright\arraybackslash}p{0.43\textwidth}@{}}\toprule
Name & Kind & Description \\\midrule\endfirsthead
\toprule Name & Kind & Description \\\midrule\endhead
\bottomrule\endfoot
\code{pca\_projection} & function & Project the data and the cluster prototypes onto the first principal components \\
\code{Sammon} & class & Sammon's nonlinear mapping to a low-dimensional space \\
\code{FuzzySammon} & class & Fuzzy Sammon mapping (FUZZSAM) of a fitted fuzzy clustering \\
\code{FuzzySOM} & class & Fuzzy self-organizing map: regularized fuzzy c-means on a grid \\
\code{grid\_laplacian} & function & Regularization matrix $\mathbf{L}=\mathbf{G}^T\mathbf{G}$ of a rectangular grid \\
\code{sammon\_stress} & function & Sammon's stress of an embedding Y of the data X \\
\code{recomputed\_membership} & function & Fuzzy memberships recomputed in the projected space \\
\code{projection\_quality} & function & Measures of how well a 2-D (or q-D) projection represents a fuzzy partition \\
\code{plot\_projection} & function & Scatter plot of a 2-D projection with projected prototypes and membership contours \\
\code{plot\_stress} & function & Stress versus iteration for one or several stress\_history\_ arrays \\
\code{plot\_map} & function & Draw a fuzzy SOM as a colored grid: U-matrix, a component plane or the hit histogram \\
\code{plot\_som\_lattice} & function & Prototypes of a fuzzy SOM in the data space, with grid neighbors connected \\
\code{make\_ellipsoids} & function & Elongated Gaussian clusters with random orientations in $n$ dimensions \\
\code{make\_spiral\_3d} & function & Noisy points along a part of a helix in 3-D (Mig\'{a}ly et al., first example) \\
\code{make\_hemisphere} & function & Noisy points on the upper half of the unit sphere (Mig\'{a}ly et al., second example) \\
\end{longtable}

\section[\code{fuzzyclust.models}]{\code{fuzzyclust.models}: Model identification}
\begin{longtable}{@{}>{\raggedright\arraybackslash}p{0.39\textwidth} l >{\raggedright\arraybackslash}p{0.43\textwidth}@{}}\toprule
Name & Kind & Description \\\midrule\endfirsthead
\toprule Name & Kind & Description \\\midrule\endhead
\bottomrule\endfoot
\code{ModifiedGG} & class & Modified Gath--Geva clustering: clusters are interpretable Takagi--Sugeno rules \\
\code{TakagiSugeno} & class & Takagi--Sugeno (TS) fuzzy model with rules extracted from fuzzy clusters \\
\code{FuzzyCRegression} & class & Fuzzy c-regression models (FCRM): clusters are regression hyperplanes \\
\code{SupervisedFuzzyClassifier} & class & Fuzzy rule-based classifier identified by supervised Gath--Geva type clustering \\
\code{FuzzyPartitionDiscretizer} & class & Supervised fuzzy discretization of continuous features \\
\code{HingingHyperplaneTree} & class & Binary regression tree with hinging-hyperplane nodes identified by fuzzy c-regression \\
\code{HingeNode} & class & One node of a HingingHyperplaneTree \\
\code{simulate\_free\_run} & function & Free-run (parallel) simulation of a NARX model \\
\code{weighted\_least\_squares} & function & Weighted least-squares estimate of local model parameters \\
\code{select\_features\_by\_separability} & function & Backward feature elimination by the Fisher interclass separability of the rules \\
\code{supervised\_fuzzy\_partition} & function & Per-feature fuzzy partitions derived from supervised fuzzy clustering \\
\code{ruspini\_partition} & function & Ruspini-type (equidistant triangular) partition of every feature, the unsupervised baseline \\
\code{membership\_1d} & function & Evaluate the fuzzy sets of one feature at the values x \\
\code{set\_similarity} & function & Jaccard-type similarity $S=|A\cap B|/|A\cup B|$ of two fuzzy sets sampled on a common grid \\
\code{plot\_membership\_functions} & function & Gaussian antecedent membership functions $A_{ij}(x_j)$ of every rule, one panel per input \\
\code{plot\_local\_models} & function & One-input Takagi--Sugeno model: data, local linear models, model output and firing strengths \\
\code{plot\_fuzzy\_partition} & function & Fuzzy sets of one feature of a fitted FuzzyPartitionDiscretizer \\
\code{plot\_hinge\_tree\_surface} & function & Prediction surface (filled contours) of a two-input hinging-hyperplane tree with its hinges \\
\code{plot\_decision\_regions} & function & Decision regions of a two-input classifier, optionally with the rule centers and antecedent ellipses \\
\code{make\_sinc} & function & Samples of the sinc function, the classic static benchmark of fuzzy modeling \\
\code{make\_piecewise\_linear} & function & A continuous piecewise-linear function of two variables (sum of two hinges) \\
\code{piecewise\_linear\_function} & function & Noise-free value of the function behind make\_piecewise\_linear \\
\code{make\_switching\_lines} & function & Switching regression data: samples of several lines over a common input range \\
\code{make\_second\_order\_plant} & function & The second-order nonlinear benchmark plant of Abonyi, Babu\v{s}ka and Szeifert (2002) \\
\code{second\_order\_plant\_function} & function & Nonlinear component $f(y(k-1), y(k-2))$ of make\_second\_order\_plant \\
\code{make\_mimo\_plant} & function & A two-input two-output nonlinear process with operating-point-dependent gains \\
\code{make\_two\_class\_overlap} & function & Two overlapping classes, each a mixture of two Gaussian blobs (XOR layout) \\
\code{make\_banded\_classes} & function & Three classes separated by thresholds on two features, plus one irrelevant feature \\
\end{longtable}

\section[\code{fuzzyclust.timeseries}]{\code{fuzzyclust.timeseries}: Time series}
\begin{longtable}{@{}>{\raggedright\arraybackslash}p{0.39\textwidth} l >{\raggedright\arraybackslash}p{0.43\textwidth}@{}}\toprule
Name & Kind & Description \\\midrule\endfirsthead
\toprule Name & Kind & Description \\\midrule\endhead
\bottomrule\endfoot
\code{FuzzySegmentation} & class & Fuzzy segmentation of a multivariate time series (modified Gath--Geva clustering) \\
\code{ChaoticPredictor} & class & Operating-regime model of a chaotic series in the reconstructed state space \\
\code{MixedInputStructureSelection} & class & Estimator form of select\_structure\_mixed \\
\code{bottom\_up\_segmentation} & function & Bottom-up segmentation with PCA models (Keogh et al. style, multivariate extension) \\
\code{pca\_segment\_cost} & function & Cost of approximating a segment by one PCA model with n\_components components \\
\code{false\_nearest\_neighbors} & function & Classical false-nearest-neighbor percentages over the model orders $(n_a, n_b)$ \\
\code{fnn\_model\_order} & function & Model order of a NARX model by false nearest neighbors with a cluster-based threshold \\
\code{select\_model\_order} & function & Convenience wrapper of fnn\_model\_order returning only the selected orders \\
\code{cluster\_based\_threshold} & function & Point-wise FNN thresholds $R_k$ from Gath--Geva clusters in the regression space \\
\code{delay\_embedding} & function & Delay vectors $\mathbf{y}_k=[y_k, y_{k-\tau},\dots,y_{k-\tau(d_e-1)}]^T$ of a scalar series \\
\code{select\_delay} & function & Lag $\tau$ of the delay embedding from the series itself \\
\code{embedding\_scan} & function & Fit ChaoticPredictor for several embedding dimensions and tabulate the criteria \\
\code{select\_embedding\_dimension} & function & Read the embedding and the local dimension off an embedding\_scan table \\
\code{correlation\_dimension} & function & Grassberger--Procaccia correlation dimension of a point cloud (e.g. delay vectors) \\
\code{select\_structure\_mixed} & function & Rank all candidate input structures by their fuzzy variance (Girimonte, Babu\v{s}ka and Abonyi) \\
\code{fuzzy\_variance} & function & Total fuzzy variance $\sigma_t$ of y over the fuzzy segments spanned by several inputs \\
\code{plot\_segmentation} & function & Time plots of the variables with the fuzzy segments underneath \\
\code{plot\_fnn\_table} & function & Heat map of the percentage of false nearest neighbors over the orders $(n_a, n_b)$ \\
\code{plot\_embedding\_scan} & function & One-step prediction error and local-dimension index versus the embedding dimension \\
\code{plot\_phase\_space} & function & Two coordinates of a reconstructed (or true) state space, colored by cluster \\
\code{plot\_structure\_ranking} & function & Variance--complexity trade-off $VCT=\sigma_t n_i$ along the ranking of structures \\
\code{make\_piecewise\_timeseries} & function & Multivariate series with planted regime changes (segmentation benchmark) \\
\code{make\_mackey\_glass} & function & Mackey--Glass delay-differential series (chaotic for $\tau=17$) \\
\code{make\_lorenz} & function & Trajectory of the Lorenz system (RK4, fixed step) \\
\code{make\_rossler} & function & Trajectory of the R\"{o}ssler system with the parameters of Abonyi et al. (2004) \\
\code{make\_nonlinear\_io\_system} & function & Continuous polymerization reactor benchmark of the FNN model-order paper \\
\code{make\_mixed\_input\_system} & function & Regression problem with one discrete and several continuous inputs \\
\end{longtable}

\section[\code{fuzzyclust.relational}]{\code{fuzzyclust.relational}: Relational, categorical, and mixed data}
\begin{longtable}{@{}>{\raggedright\arraybackslash}p{0.39\textwidth} l >{\raggedright\arraybackslash}p{0.43\textwidth}@{}}\toprule
Name & Kind & Description \\\midrule\endfirsthead
\toprule Name & Kind & Description \\\midrule\endhead
\bottomrule\endfoot
\code{RelationalFCM} & class & Relational fuzzy c-means (Hathaway, Davenport and Bezdek, 1989) \\
\code{NERFCM} & class & Non-Euclidean relational fuzzy c-means (Hathaway and Bezdek, 1994) \\
\code{FuzzyCMedoids} & class & Fuzzy c-medoids clustering of relational data (Krishnapuram et al., 2001) \\
\code{FuzzyKModes} & class & Fuzzy k-modes clustering of categorical data (Huang and Ng, 1999) \\
\code{FuzzyKPrototypes} & class & Fuzzy k-prototypes clustering of mixed numeric and categorical data (Huang, 1998) \\
\code{relational\_distances} & function & Squared relational distances of the training objects to relational prototypes \\
\code{matching\_dissimilarity} & function & Simple matching dissimilarity $\delta(\mathbf{a},\mathbf{b})=\sum_j[a_j\neq b_j]$ between rows \\
\code{pairwise\_dissimilarity} & function & Pairwise dissimilarities of feature vectors by any SciPy metric \\
\code{check\_dissimilarity\_matrix} & function & Validate a (square or rectangular) dissimilarity matrix \\
\code{dtw\_distance\_matrix} & function & Dynamic time warping distances between all pairs of (equal-length) series \\
\code{dtw\_alignment} & function & Optimal warping path between two series (for inspection and plotting) \\
\code{autocorrelation\_features} & function & Sample autocorrelations $\hat\rho(1),\dots,\hat\rho(L)$ of every series \\
\code{autocorrelation\_distance\_matrix} & function & Autocorrelation-based distance between time series (D'Urso and Maharaj, 2009) \\
\code{gower\_distance\_matrix} & function & Gower's general coefficient of dissimilarity for mixed data (Gower, 1971) \\
\code{plot\_dissimilarity\_matrix} & function & Heat map of a dissimilarity matrix with the objects ordered by their labels \\
\code{plot\_dtw\_alignment} & function & Two series with the dynamic-time-warping correspondences between them \\
\code{plot\_modes} & function & Table of the fitted modes, shaded by how much membership mass the mode carries \\
\code{plot\_series\_clusters} & function & Time series drawn in the color of their cluster, medoid series emphasized \\
\code{make\_categorical\_blobs} & function & Categorical records scattered around planted modes \\
\code{make\_mixed\_blobs} & function & Records with Gaussian numeric attributes and categorical attributes around modes \\
\code{make\_warped\_series} & function & Time series of three shape classes with random time warps \\
\end{longtable}

\section[\code{fuzzyclust.incomplete}]{\code{fuzzyclust.incomplete}: Incomplete data}
\begin{longtable}{@{}>{\raggedright\arraybackslash}p{0.39\textwidth} l >{\raggedright\arraybackslash}p{0.43\textwidth}@{}}\toprule
Name & Kind & Description \\\midrule\endfirsthead
\toprule Name & Kind & Description \\\midrule\endhead
\bottomrule\endfoot
\code{IncompleteFCM} & class & Fuzzy c-means clustering of incomplete data (Hathaway and Bezdek, 2001) \\
\code{STRATEGIES} & constant & Names of the missing-data strategies of \code{IncompleteFCM} (WDS, PDS, OCS, NPS) \\
\code{partial\_distances} & function & Squared \emph{partial distances} of (incomplete) samples to prototypes \\
\code{make\_mcar} & function & Remove values \emph{missing completely at random} (MCAR) from a data matrix \\
\code{make\_incomplete\_blobs} & function & Isotropic Gaussian clusters with values missing completely at random \\
\code{plot\_incomplete\_clusters} & function & Scatter plot of a clustering of incomplete 2-D data at the completed coordinates \\
\code{plot\_imputation} & function & Imputed against true values of the missing entries, one color per feature \\
\end{longtable}

\section[\code{fuzzyclust.large}]{\code{fuzzyclust.large}: Very large data and streams}
\begin{longtable}{@{}>{\raggedright\arraybackslash}p{0.39\textwidth} l >{\raggedright\arraybackslash}p{0.43\textwidth}@{}}\toprule
Name & Kind & Description \\\midrule\endfirsthead
\toprule Name & Kind & Description \\\midrule\endhead
\bottomrule\endfoot
\code{WeightedFCM} & class & Weighted fuzzy c-means (wFCM): fuzzy c-means with a weight per sample \\
\code{RandomSampleFCM} & class & Random sample and extend FCM (rseFCM) for very large data \\
\code{SinglePassFCM} & class & Single-pass fuzzy c-means (spFCM) for very large and streaming data \\
\code{OnlineFCM} & class & Online fuzzy c-means (oFCM) for very large and streaming data \\
\code{extend\_partition} & function & Memberships of (many) samples in fixed prototypes (the \emph{extension} step) \\
\code{make\_large\_blobs} & function & Many samples from a few isotropic Gaussian clusters, in random order \\
\code{make\_drifting\_stream} & function & A stream of chunks whose cluster centers drift from chunk to chunk \\
\code{plot\_chunk\_timeline} & function & Prototype coordinate against the chunk index for a chunked model \\
\code{plot\_center\_trajectories} & function & Trajectories of the prototypes over the chunks on top of (a subsample of) the 2-D data \\
\end{longtable}

\section[\code{fuzzyclust.image}]{\code{fuzzyclust.image}: Image segmentation}
\begin{longtable}{@{}>{\raggedright\arraybackslash}p{0.39\textwidth} l >{\raggedright\arraybackslash}p{0.43\textwidth}@{}}\toprule
Name & Kind & Description \\\midrule\endfirsthead
\toprule Name & Kind & Description \\\midrule\endhead
\bottomrule\endfoot
\code{SpatialFCM} & class & Spatially constrained fuzzy c-means for image segmentation \\
\code{METHODS} & constant & Names of the spatial methods of \code{SpatialFCM} \\
\code{neighborhood\_sum} & function & Sum of an image over the neighbors of every pixel, $\sum_{r\in N_k} y_r$ (center excluded) \\
\code{weighted\_sum\_image} & function & Linearly weighted sum image of EnFCM \\
\code{local\_similarity\_image} & function & Locally adaptive weighted image of FGFCM, $\xi_k=\sum_{r\in N_k}S_{kr}x_r/\sum_{r}S_{kr}$ \\
\code{make\_synthetic\_image} & function & A piecewise-constant image of simple shapes with known labels, plus a noisy copy \\
\code{add\_image\_noise} & function & Corrupt an image with Gaussian, salt-and-pepper or both kinds of noise \\
\code{segmentation\_accuracy} & function & Fraction of correctly labeled pixels after the best matching of the labels \\
\code{match\_labels} & function & Best one-to-one mapping of predicted labels onto true labels (Hungarian algorithm) \\
\code{plot\_segmentation} & function & Panels original / noisy / labels / membership maps of a fitted SpatialFCM \\
\code{plot\_image} & function & Show a gray-level or RGB image without axes ticks \\
\code{plot\_label\_image} & function & Show a label image with the fixed cluster palette (cluster \emph{i} in color \emph{i}) \\
\code{plot\_membership\_maps} & function & One panel per cluster with its membership image in [0, 1] \\
\end{longtable}

\section[\code{fuzzyclust.semisupervised}]{\code{fuzzyclust.semisupervised}: Partial supervision}
\begin{longtable}{@{}>{\raggedright\arraybackslash}p{0.39\textwidth} l >{\raggedright\arraybackslash}p{0.43\textwidth}@{}}\toprule
Name & Kind & Description \\\midrule\endfirsthead
\toprule Name & Kind & Description \\\midrule\endhead
\bottomrule\endfoot
\code{SemiSupervisedFCM} & class & Fuzzy c-means with partial supervision (Pedrycz and Waletzky, 1997) \\
\code{PairwiseConstrainedFCM} & class & Fuzzy clustering with pairwise constraints, PCCA (Grira, Crucianu, and Boujemaa) \\
\code{check\_pairwise\_constraints} & function & Validate lists of pairwise constraints and return them as integer arrays (n\_pairs, 2) \\
\code{constraint\_satisfaction} & function & Fraction of satisfied must-link and cannot-link pairs under hard labels \\
\code{make\_partially\_labeled\_pair} & function & Two long parallel groups whose \emph{natural} split disagrees with the classes \\
\code{make\_partial\_labels} & function & Hide most labels of y: keep a random subset, set the rest to -1 \\
\code{make\_pairwise\_constraints} & function & Draw random must-link and cannot-link pairs consistent with the labels y \\
\code{plot\_label\_propagation} & function & Labeled seeds as large black-edged markers, the rest colored by the cluster they were pulled into \\
\code{plot\_pairwise\_constraints} & function & Overlay must-link (solid green) and cannot-link (dashed vermilion) pairs on a scatter plot \\
\end{longtable}

\section[\code{fuzzyclust.coclustering}]{\code{fuzzyclust.coclustering}: Co-clustering}
\begin{longtable}{@{}>{\raggedright\arraybackslash}p{0.39\textwidth} l >{\raggedright\arraybackslash}p{0.43\textwidth}@{}}\toprule
Name & Kind & Description \\\midrule\endfirsthead
\toprule Name & Kind & Description \\\midrule\endhead
\bottomrule\endfoot
\code{FCCM} & class & Fuzzy clustering for categorical multivariate data, FCCM (Oh, Honda, and Ichihashi, 2001) \\
\code{FCoDoK} & class & Fuzzy co-clustering of documents and keywords, Fuzzy CoDoK (Kummamuru, Dhawale, Krishnapuram, 2003) \\
\code{make\_block\_matrix} & function & A block-structured object--item count matrix with known co-clusters \\
\code{plot\_coclusters} & function & Heat map of an object--item matrix with rows and columns reordered by co-cluster \\
\code{plot\_item\_memberships} & function & Bar chart of the item memberships of every co-cluster \\
\end{longtable}

\section[\code{fuzzyclust.collaborative}]{\code{fuzzyclust.collaborative}: Collaborative clustering}
\begin{longtable}{@{}>{\raggedright\arraybackslash}p{0.39\textwidth} l >{\raggedright\arraybackslash}p{0.43\textwidth}@{}}\toprule
Name & Kind & Description \\\midrule\endfirsthead
\toprule Name & Kind & Description \\\midrule\endhead
\bottomrule\endfoot
\code{CollaborativeFCM} & class & Horizontal collaborative fuzzy clustering (Pedrycz, 2002) \\
\code{align\_partitions} & function & Permute the clusters of U to match those of U\_ref (maximal membership overlap) \\
\code{make\_collaborative\_sites} & function & The same objects described at several data sites by different features \\
\code{plot\_site\_partitions} & function & Scatter plot of every site's (two-dimensional) data colored by its own partition \\
\code{plot\_collaboration\_trajectory} & function & Agreement between the sites as a function of the collaboration weight $\alpha$ \\
\end{longtable}

\section[\code{fuzzyclust.ensemble}]{\code{fuzzyclust.ensemble}: Cluster ensembles}
\begin{longtable}{@{}>{\raggedright\arraybackslash}p{0.39\textwidth} l >{\raggedright\arraybackslash}p{0.43\textwidth}@{}}\toprule
Name & Kind & Description \\\midrule\endfirsthead
\toprule Name & Kind & Description \\\midrule\endhead
\bottomrule\endfoot
\code{FuzzyClusterEnsemble} & class & Consensus of a library of soft clusterings (Punera and Ghosh, 2008) \\
\code{coassociation\_matrix} & function & Soft co-association (similarity) matrix of a library of fuzzy partitions \\
\code{evidence\_accumulation} & function & Evidence accumulation clustering of a co-association matrix (Fred and Jain, 2005) \\
\code{soft\_mcla} & function & Soft meta-clustering algorithm, sMCLA (Punera and Ghosh, 2008) \\
\code{information\_theoretic\_kmeans} & function & Information-theoretic k-means consensus, ITK (Punera and Ghosh, 2008) \\
\code{cspa} & function & Cluster-based similarity partitioning, CSPA (Strehl and Ghosh, 2002), for hard partitions \\
\code{mcla} & function & Meta-clustering algorithm, MCLA (Strehl and Ghosh, 2002), for hard partitions \\
\code{stability\_curve} & function & Reproducibility of the consensus partition as the library grows \\
\code{make\_overlapping\_blobs} & function & Gaussian groups that overlap enough for single clustering runs to disagree \\
\code{plot\_coassociation} & function & Heat map of a co-association matrix with the samples reordered by (consensus) label \\
\code{plot\_lifetimes} & function & Bar chart of the lifetime of every number of clusters in the co-association dendrogram \\
\end{longtable}

\section[\code{fuzzyclust.datasets}]{\code{fuzzyclust.datasets}: Data sets}
\begin{longtable}{@{}>{\raggedright\arraybackslash}p{0.39\textwidth} l >{\raggedright\arraybackslash}p{0.43\textwidth}@{}}\toprule
Name & Kind & Description \\\midrule\endfirsthead
\toprule Name & Kind & Description \\\midrule\endhead
\bottomrule\endfoot
\code{make\_fuzzy\_blobs} & function & Isotropic Gaussian clusters with random centers (MATLAB nDexample) \\
\code{make\_ellipses} & function & Elongated 2-D Gaussian clusters with different orientations \\
\code{make\_different\_densities} & function & Clusters of very different size (density) and spread \\
\code{make\_rings} & function & Concentric rings (hollow clusters) in 2-D \\
\code{make\_lines} & function & Noisy line segments in 2-D (linear sub-space prototypes) \\
\code{make\_noisy\_blobs} & function & Gaussian clusters plus uniformly distributed outliers \\
\code{make\_curved\_manifold} & function & Points along two interleaved spiral arms (nonconvex clusters) \\
\code{make\_two\_moons} & function & Two interleaving half circles (thin wrapper around scikit-learn) \\
\code{load\_motorcycle} & function & Motorcycle impact data (Silverman, 1985) used throughout the MATLAB toolbox \\
\code{load\_iris} & function & Fisher's iris data ($150\times 4$, 3 classes) via scikit-learn \\
\code{load\_wine} & function & Wine recognition data ($178\times 13$, 3 classes) via scikit-learn \\
\code{load\_breast\_cancer} & function & Wisconsin breast cancer data ($569\times 30$, 2 classes) via scikit-learn \\
\end{longtable}

\section[\code{fuzzyclust.preprocessing}]{\code{fuzzyclust.preprocessing}: Preprocessing}
\begin{longtable}{@{}>{\raggedright\arraybackslash}p{0.39\textwidth} l >{\raggedright\arraybackslash}p{0.43\textwidth}@{}}\toprule
Name & Kind & Description \\\midrule\endfirsthead
\toprule Name & Kind & Description \\\midrule\endhead
\bottomrule\endfoot
\code{normalize} & function & Scale the data and return the scaler for inverse transformation \\
\code{make\_regressors} & function & Regression matrix and target of a NARX model \\
\code{lagged\_matrix} & function & Delay-embedding of a univariate series \\
\end{longtable}

\section[\code{fuzzyclust.plotting}]{\code{fuzzyclust.plotting}: Plotting helpers}
\begin{longtable}{@{}>{\raggedright\arraybackslash}p{0.39\textwidth} l >{\raggedright\arraybackslash}p{0.43\textwidth}@{}}\toprule
Name & Kind & Description \\\midrule\endfirsthead
\toprule Name & Kind & Description \\\midrule\endhead
\bottomrule\endfoot
\code{PALETTE} & constant & Fixed categorical color palette; cluster $i$ always gets \code{PALETTE[i]} \\
\code{plot\_clusters} & function & Scatter plot of 2-D data colored by (hard) cluster label with prototypes \\
\code{plot\_membership\_contours} & function & Contour map of the maximal membership $\max_i \mu_i(\mathbf{x})$ over the plane \\
\code{plot\_covariance\_ellipses} & function & Draw the $n_\text{std}$-sigma ellipses of (fuzzy) covariance matrices \\
\code{plot\_validity\_scan} & function & Plot validity indices versus the number of clusters \\
\code{plot\_objective} & function & Objective $J$ versus iteration for one or several fitted estimators \\
\code{plot\_membership\_profiles} & function & Membership degrees $\mu_{ik}$ of every cluster along the sample index \\
\end{longtable}
}
